\section{视觉与资源管理}
\subsection{关键画面与时间戳}
本讲依赖封面图与 4 张关键帧来加强视觉证据。在没有 slides 的情况下，这些帧展示了 self-rewarding、self-instruct、meta-judge 等核心概念。附表列出帧名、时间区间与文本说明，方便后续追溯。

\begin{table}[H]
\centering
\begin{tabular}{p{3cm}p{3cm}p{7cm}}
\toprule
画面文件 & 时间 & 说明 \\
\midrule
frame-01-selfreward.jpg & 00:02:35--00:02:37 & self-improvement pipeline 的概念图 \\
frame-02-iterative.jpg & 00:36:00--00:36:02 & Self-Rewarding 迭代闭环 \\
frame-03-selfinstr.jpg & 00:44:40--00:44:42 & Self-Instruct 任务与 verified response \\
frame-04-meta.jpg & 01:01:30--01:01:32 & Meta Judge 对比 draft 与 verified response \\
\bottomrule
\end{tabular}
\caption{本讲使用的关键帧与时间标注。}
\end{table}

\subsection{字幕与写作支持}
我们用 \texttt{lecture11.en.srt} 清洗出 \texttt{subs\_clean.txt}，以 2500 行/块为单位提炼主题句，再按照教学逻辑进行章节划分。clean 版文本减少重复（如 ``days I'm using AI...''），为扩充每节内容提供了稳定素材。

\subsection{本章小结}
视觉资产与字幕清洗构成了补充素材的双重支撑：前者满足用户要求的 frame + 时间；后者让章节写作有据可依，保持技术细节与顺序的一致性。

